% ============================================================================
% LEHRERHINWEISE: Las asignaturas (Schulfächer)
% ============================================================================
% Sequenz 4: Mi colegio | Block a
% Fachfarbe: #c41e3a (Karminrot)
% ============================================================================

\documentclass[11pt, a4paper]{article}

% ============================================================================
% SW-MODUS TOGGLE
% ============================================================================
\newif\ifswmodus
\swmodusfalse

% ============================================================================
% PAKETE
% ============================================================================

\usepackage[utf8]{inputenc}
\usepackage[T1]{fontenc}
\usepackage[ngerman]{babel}
\usepackage{helvet}
\renewcommand{\familydefault}{\sfdefault}

\usepackage[margin=2cm]{geometry}
\usepackage{enumitem}
\usepackage{tabularx}
\usepackage{booktabs}
\usepackage{longtable}

\usepackage{xcolor}
\usepackage{amssymb}
\usepackage{tcolorbox}
\tcbuselibrary{skins}

\usepackage{fancyhdr}
\usepackage{lastpage}
\usepackage{needspace}

% ============================================================================
% FARBEN
% ============================================================================

\definecolor{spanisch-color}{HTML}{c41e3a}
\definecolor{grau-color}{HTML}{888888}
\definecolor{hellgrau-color}{HTML}{f5f5f5}
\definecolor{basis-color}{HTML}{2a7a4b}
\definecolor{standard-color}{HTML}{2c5aa0}
\definecolor{erweiterung-color}{HTML}{7b1fa2}
\definecolor{loesung-color}{HTML}{c62828}

\definecolor{sw-dark}{HTML}{000000}
\definecolor{sw-gray}{HTML}{666666}
\definecolor{sw-white}{HTML}{ffffff}

\ifswmodus
    \colorlet{spanisch}{sw-dark}
    \colorlet{grau}{sw-gray}
    \colorlet{hellgrau}{sw-white}
    \colorlet{basis}{sw-dark}
    \colorlet{standard}{sw-dark}
    \colorlet{erweiterung}{sw-dark}
    \colorlet{loesung}{sw-dark}
\else
    \colorlet{spanisch}{spanisch-color}
    \colorlet{grau}{grau-color}
    \colorlet{hellgrau}{hellgrau-color}
    \colorlet{basis}{basis-color}
    \colorlet{standard}{standard-color}
    \colorlet{erweiterung}{erweiterung-color}
    \colorlet{loesung}{loesung-color}
\fi

\newcommand{\fachfarbe}{spanisch}

% ============================================================================
% VARIABLEN
% ============================================================================

\newcommand{\thema}{Las asignaturas (Schulfächer)}
\newcommand{\sequenz}{04}
\newcommand{\block}{a}
\newcommand{\fach}{Spanisch}
\newcommand{\klasse}{7}
\newcommand{\dauer}{90 Min. (Doppelstunde)}
\newcommand{\lerngruppengroesse}{24}
\newcommand{\datum}{\today}

% ============================================================================
% KOPF- UND FUSSZEILE
% ============================================================================

\pagestyle{fancy}
\fancyhf{}
\renewcommand{\headrulewidth}{0.4pt}
\renewcommand{\footrulewidth}{0.4pt}

\lhead{\fach{} -- Klasse \klasse}
\chead{\textbf{LH-\sequenz\block{} (Lehrerhinweise)}}
\rhead{\datum}

\lfoot{}
\cfoot{}
\rfoot{Seite \thepage\ von \pageref{LastPage}}

% ============================================================================
% BEFEHLE
% ============================================================================

\newcommand{\B}{\textcolor{basis}{\textbf{[B]}}}
\newcommand{\Std}{\textcolor{standard}{\textbf{[S]}}}
\newcommand{\E}{\textcolor{erweiterung}{\textbf{[E]}}}

\newcommand{\loesung}[1]{\textcolor{loesung}{\textbf{#1}}}

\newcommand{\es}[1]{\textcolor{\fachfarbe}{\textbf{#1}}}

\newtcolorbox{kurzplan}{
    colback=hellgrau,
    colframe=\fachfarbe,
    colbacktitle=white,
    coltitle=\fachfarbe,
    boxrule=1pt,
    arc=0pt,
    fonttitle=\bfseries,
    attach boxed title to top left={yshift=-2mm, xshift=5mm},
    boxed title style={boxrule=0pt, colback=white},
    title=Kurzplan
}

\newtcolorbox{differenzierung}{
    colback=white,
    colframe=grau,
    colbacktitle=white,
    coltitle=grau,
    boxrule=0.5pt,
    arc=0pt,
    fonttitle=\bfseries,
    attach boxed title to top left={yshift=-2mm, xshift=5mm},
    boxed title style={boxrule=0pt, colback=white},
    title=Differenzierung
}

\newtcolorbox{fehlerbox}{
    colback=white,
    colframe=loesung,
    colbacktitle=white,
    coltitle=loesung,
    boxrule=0.5pt,
    arc=0pt,
    fonttitle=\bfseries,
    attach boxed title to top left={yshift=-2mm, xshift=5mm},
    boxed title style={boxrule=0pt, colback=white},
    title=Häufige Fehler
}

\newtcolorbox{materialbox}{
    colback=white,
    colframe=\fachfarbe,
    colbacktitle=white,
    coltitle=\fachfarbe,
    boxrule=0.5pt,
    arc=0pt,
    fonttitle=\bfseries,
    attach boxed title to top left={yshift=-2mm, xshift=5mm},
    boxed title style={boxrule=0pt, colback=white},
    title=Material-Checkliste
}

% ============================================================================
% DOKUMENT
% ============================================================================

\begin{document}

% ============================================================================
% KOPFBEREICH
% ============================================================================

\begin{center}
    {\Large\bfseries\textcolor{\fachfarbe}{Lehrerhinweise: \thema}}\\[0.3em]
    {\normalsize LH-\sequenz\block{} | \dauer}
\end{center}

\vspace{10pt}

% ============================================================================
% 1. KURZPLAN
% ============================================================================

\section*{1. Stundenverlauf (Kurzplan)}

\textbf{1. Stunde (45 Min.)}

\begin{tabularx}{\textwidth}{|c|l|X|l|}
\hline
\textbf{Zeit} & \textbf{Phase} & \textbf{Inhalt} & \textbf{Material} \\
\hline
0--5 & Einstieg & Bildimpuls Stundenplan, Thema einführen & Tafel, Bild \\
\hline
5--15 & Input & Schulfächer mit Bildkarten präsentieren, nachsprechen & Bildkarten \\
\hline
15--25 & Übung 1 & Memory-Spiel in 4er-Gruppen & Memory \\
\hline
25--35 & Übung 2 & AB Aufgabe 1 (Zuordnung) & AB-04a \\
\hline
35--40 & Sicherung & Lösungen besprechen, Aussprache üben & Tafel \\
\hline
40--45 & Überleitung & Artikel sammeln, Ausblick: Meinungen & Tafel \\
\hline
\end{tabularx}

\vspace{1em}

\textbf{2. Stunde (45 Min.)}

\begin{tabularx}{\textwidth}{|c|l|X|l|}
\hline
\textbf{Zeit} & \textbf{Phase} & \textbf{Inhalt} & \textbf{Material} \\
\hline
0--5 & Wiederholung & Schnelles Vokabelquiz mit Bildkarten & Bildkarten \\
\hline
5--15 & Input & ``Me gusta(n)'' einführen, Beispiele & Tafel \\
\hline
15--25 & Übung 3 & AB Aufgabe 3 (Meinung schreiben) & AB-04a \\
\hline
25--35 & Anwendung & Partnerarbeit: Mini-Dialoge über Fächer & -- \\
\hline
35--42 & Festigung & AB Aufgabe 4 (Stundenplan lesen) & AB-04a \\
\hline
42--45 & Abschluss & Zusammenfassung, HA ansagen & -- \\
\hline
\end{tabularx}

% ============================================================================
% 2. DIFFERENZIERUNG
% ============================================================================

\section*{2. Differenzierung}

\begin{differenzierung}
\textbf{Legende:} \B{} Basis (BR) | \Std{} Standard (MR) | \E{} Erweiterung (GY)

\vspace{0.5em}

\textbf{WICHTIG:} Diese Marker sind NUR für die Lehrkraft. Auf Schülermaterial erscheinen KEINE Niveaumarkierungen!
\end{differenzierung}

\vspace{1em}

\begin{tabularx}{\textwidth}{|l|X|X|X|}
\hline
& \B{} \textbf{Basis} & \Std{} \textbf{Standard} & \E{} \textbf{Erweiterung} \\
\hline
\textbf{Aufgabe 1} & Mit Bildkarten als Hilfe & Selbstständig & + Sätze bilden \\
\hline
\textbf{Aufgabe 2} & Partnerhilfe erlaubt & Selbstständig & + Erklärung maskulin/feminin \\
\hline
\textbf{Aufgabe 3} & Vorgegebene Fächer & Freie Wahl & + Begründung (porque...) \\
\hline
\textbf{Aufgabe 4} & a) und b) & a), b), c) & + eigener Stundenplan \\
\hline
\end{tabularx}

\vspace{1em}

\textbf{Konkrete Hinweise:}
\begin{itemize}
    \item \B{} SuS mit LRS: Zeitzugabe (+10 Min.), vergrößerte AB-Version (A3)
    \item \B{} DaZ-SuS: Vorab deutsche Fächernamen klären, Bildwörterbuch nutzen
    \item \E{} Leistungsstarke: Begründungen einfordern: ``Me gusta la musica \textbf{porque es interesante}.''
    \item \E{} Zusatz: Eigenen Stundenplan auf Spanisch erstellen
\end{itemize}

% ============================================================================
% 3. HÄUFIGE FEHLER
% ============================================================================

\section*{3. Häufige Fehler}

\begin{fehlerbox}
\begin{tabularx}{\textwidth}{|l|X|X|}
\hline
\textbf{Fehler} & \textbf{Beispiel} & \textbf{Reaktion} \\
\hline
Betonung matematicas & *maTEmaticas statt \es{mateMÁticas} & Silben klatschen: ma-te-MÁ-ti-cas \\
\hline
gusta vs. gustan & *Me gusta las matematicas & Visualisierung: 1 Fach = gusta, viele = gustan \\
\hline
Artikel vergessen & *Me gusta musica & Immer mit Artikel lernen: ``la musica'' \\
\hline
Interferenz ``ich mag'' & *Yo gusto espanol & Bewusstmachung: ``Das Fach gefällt mir'' \\
\hline
educacion fisica & *educacion physica & Hinweis: ``ph'' gibt es im Spanischen nicht \\
\hline
\end{tabularx}
\end{fehlerbox}

% ============================================================================
% 4. MATERIAL-CHECKLISTE
% ============================================================================

\section*{4. Material-Checkliste}

\begin{materialbox}
\begin{itemize}[label=$\square$]
    \item AB-\sequenz\block.pdf (\lerngruppengroesse{} Kopien)
    \item ML-\sequenz\block.pdf (1$\times$ für Lehrkraft)
    \item Lehrwerk: Qué pasa 1, S. 36-42
    \item Audio-Track 18 (optional)
    \item Bildkarten Schulfächer (11 Stück)
    \item Memory-Karten (6 Sets für Gruppenarbeit)
    \item Tafelmarker / Kreide
    \item Optional: LP-\sequenz\block.html (Tablets/PCs)
\end{itemize}
\end{materialbox}

% ============================================================================
% 5. LÖSUNGEN
% ============================================================================

\section*{5. Lösungen AB-\sequenz\block}

\textbf{Merkkasten: Schulfächer}

\begin{tabular}{ll}
\es{las matematicas} & = Mathematik \\
\es{el espanol} & = Spanisch \\
\es{el ingles} & = Englisch \\
\es{el aleman} & = Deutsch \\
\es{el arte} & = Kunst \\
\es{la informatica} & = Informatik \\
\es{la historia} & = Geschichte \\
\es{la geografia} & = Erdkunde \\
\es{la biologia} & = Biologie \\
\es{la musica} & = Musik \\
\es{la educacion fisica} & = Sport \\
\end{tabular}

\vspace{1em}

\textbf{Aufgabe 1: Zuordnung} (4 Punkte)

\begin{tabular}{ll}
\es{las matematicas} & $\rightarrow$ \loesung{Mathematik} \\
\es{la historia} & $\rightarrow$ \loesung{Geschichte} \\
\es{el ingles} & $\rightarrow$ \loesung{Englisch} \\
\es{la educacion fisica} & $\rightarrow$ \loesung{Sport} \\
\end{tabular}

\vspace{1em}

\textbf{Aufgabe 2: Artikel} (4 Punkte)

\begin{tabular}{llll}
a) \loesung{el} espanol & b) \loesung{la} musica & c) \loesung{la} geografia & d) \loesung{la} informatica \\
e) \loesung{la} biologia & f) \loesung{el} arte & g) \loesung{el} aleman & h) \loesung{el} ingles \\
\end{tabular}

\vspace{1em}

\textbf{Aufgabe 3: Meinung äußern} (6 Punkte)

Mögliche Antworten:
\begin{itemize}
    \item a) \loesung{Me gusta la musica. Me gusta el espanol.}
    \item b) \loesung{No me gusta la historia. No me gustan las matematicas.}
    \item c) \loesung{?Te gusta el arte? -- Si, me gusta. / No, no me gusta.}
\end{itemize}

\textit{Bewertung: Je 1P pro korrektem Satz mit richtigem Artikel und korrekter gusta(n)-Form}

\vspace{1em}

\textbf{Aufgabe 4: Stundenplan} (6 Punkte)

\begin{itemize}
    \item a) Maria tiene \loesung{(las) matematicas}. (2P)
    \item b) Maria tiene \loesung{(la) educacion fisica}. (2P)
    \item c) Maria tiene \loesung{cinco} asignaturas. (2P)
\end{itemize}

\textit{Hinweis: Bei c) auch ``5'' akzeptieren (1P), ``cinco'' ausgeschrieben (2P)}

% ============================================================================
% 6. HAUSAUFGABE
% ============================================================================

\section*{6. Hausaufgabe}

\begin{itemize}
    \item Vokabeln lernen: 11 Schulfächer mit Artikel
    \item Lehrwerk S. 38, Aufgabe 3 (schriftlich)
    \item Optional: Lernpfad LP-\sequenz\block.html (digital)
\end{itemize}

% ============================================================================
% 7. REFLEXIONSNOTIZEN
% ============================================================================

\section*{7. Reflexionsnotizen (nach Durchführung)}

\begin{tabularx}{\textwidth}{|l|X|}
\hline
\textbf{Was lief gut?} & \\[2em]
\hline
\textbf{Was lief nicht?} & \\[2em]
\hline
\textbf{Zeitmanagement} & \\[2em]
\hline
\textbf{Für nächstes Mal} & \\[2em]
\hline
\end{tabularx}

\end{document}
